\subsection{Execution and Verification Backbone}

The implementation combines three off-the-shelf components: the QCP symbolic execution engine~\cite{wu2025qcppracticalseparationlogicbased,vst-a}, an LLM, and the Frama-C/WP automated verifier~\cite{cuoq2012frama}, and orchestrates them through a thin translation layer. Symbolic execution drives the pipeline and pauses at each loop and function call, where the LLM proposes the missing annotations from the exposed symbolic state and Frama-C/WP discharges the resulting proof obligations.


The choice of QCP and Frama-C is motivated by their complementary strengths. QCP is used as the symbolic execution engine because it produces Coq-verified postconditions and, more importantly, exposes its symbolic state (path conditions plus symbolic stores) at arbitrary program points in a logical form that is directly reusable for specification synthesis. Crucially, QCP satisfies the symbolic-execution conditions of Proposition~\ref{thm:sp-sound-complete}, thereby providing the foundation for the strongest-postcondition guarantee on loop-free segments. Other widely used alternatives such as Klee~\cite{DBLP:conf/osdi/CadarDE08} are not built on Hoare logic and accordingly neither expose intermediate states at arbitrary locations nor present them in a logic-oriented form suitable for specification synthesis. QCP also has strong support for programs that manipulate complex data structures and ships with a library of verified inductive predicates and lemmas, which matches our goal of producing formally precise specifications. The gap QCP leaves --- it requires user-supplied loop invariants and callee pre-/postconditions, and it may emit lemmas that need additional Coq interactive proofs --- is exactly the gap filled by the other two components. The LLM proposes the missing invariants and specifications from the symbolic state QCP exposes; Frama-C/WP then discharges the resulting proof obligations automatically through SMT under the LLM-friendly ACSL surface syntax. The three components therefore split the work cleanly: QCP supplies the symbolic semantics, the LLM supplies the specification candidates, and Frama-C/WP closes the verification loop.

%Stitching these three components together cleanly requires a translation layer, because each of them speaks a different annotation language: QCP consumes and emits QCPL (its native annotation language), the LLM is prompted in ACSL plus natural language, and Frama-C/WP verifies ACSL. The translator maintains a bidirectional bridge between ACSL and QCPL --- forward-translating ACSL clauses into QCPL during the $\mathrm{LoopInvGen}$ phase so that QCP can use them consistently across loop and function boundaries, and back-translating QCP's symbolic-state output into ACSL skeletons (for the LLM, paired with natural-language guidance) and into ACSL clauses (for Frama-C/WP verification). A deliberate design choice is that QCPL is never shown to the LLM: the model only ever sees ACSL plus natural-language guidance derived from the symbolic state. This avoids the well-documented failure mode of an LLM confusing two formally similar but semantically distinct annotation languages, and the skeleton-plus-prose discipline doubles as a pre-structured frame so that structural well-formedness is the translator's responsibility and the model is left to fill in the logical content of each placeholder. When verification fails, error information is translated into natural language together with its original context and routed back into the next LLM call, closing the loop without ever exposing the LLM to the verifier's internal representation either.

Stitching these three components together requires a translation layer, because QCP, the LLM, and Frama‑C/WP each speak a different annotation language: QCPL (native to QCP), ACSL (native to Frama‑C/WP), and natural-language prompts (for the LLM). The translator maintains a bidirectional bridge between ACSL and QCPL: it translates ACSL clauses into QCPL so that QCP can use them across loop and function boundaries, and it converts QCP’s symbolic‑state output into ACSL skeletons (paired with natural‑language guidance for the LLM) and ACSL clauses (for Frama‑C/WP verification). A deliberate design choice is that the LLM never sees QCPL; it only operates on ACSL plus natural‑language guidance derived from the symbolic state. This avoids the common failure mode of confusing two formally similar but semantically distinct languages, and the skeleton‑plus‑prose discipline ensures structural well‑formedness before the LLM fills in the logical content. When verification fails, error information is translated back into natural language, together with its context, and routed into the next LLM call—closing the loop without exposing the verifier’s internal representation to the model either.
